%% -----------------------------------------------------------
%% Proposta curta de TCC baseada no projeto Alvesting

\documentclass[portuguese]{sbcreviews-2025}
\usepackage[misc,geometry]{ifsym}
\usepackage{aas_macros}
\usepackage[bottom]{footmisc}
\usepackage{tabularray}
\usepackage{afterpage}
\usepackage{url}
\usepackage{pifont}

\setcitestyle{square}

\definecolor{engtitle}{rgb}{0.5,0.5,0.5}
\definecolor{darkblue}{rgb}{0.0,0.0,0.0}
\hypersetup{colorlinks,breaklinks,
            linkcolor=darkblue,urlcolor=darkblue,
            anchorcolor=darkblue,citecolor=darkblue}

\jid{ROCS}
\jtitle{SBC Computing Reviews, 2026, XX:1}
\issn{2966-3938}
\doi{10.5753/rocs.2026.XXXXXX}
\copyrightstatement{This work is licensed under a Creative Commons Attribution 4.0 International License}
\jyear{2026}

\category{Proposta de Trabalho de Conclusão de Curso}
\title[Alvesting: proposta de TCC]{Alvesting: desenvolvimento de uma plataforma de competição entre modelos de linguagem aplicada ao mercado financeiro brasileiro}
\engtitle{\textcolor{engtitle}{Alvesting: development of a large language model competition platform applied to the Brazilian financial market}}

%% Campos editoriais que ainda podem exigir ajuste:
%% - volume/numero em \jtitle
%% - identificador final em \doi

\author[Zanella 2026]{
\affil{\textbf{Artur Felippe Zanella}~[Centro Universitário - Católica de Santa Catarina~|\href{mailto:artur.zanella@catolicasc.edu.br}{~{\textit{artur.zanella@catolicasc.edu.br}}}~]}
}

\begin{document}

\begin{frontmatter}

\maketitle

\begin{mail}
Engenharia de Software, Centro Universitário - Católica de Santa Catarina, Jaraguá do Sul, Santa Catarina, Brasil.
\end{mail}

\begin{abstract-pt}
Este documento apresenta a proposta de Trabalho de Conclusão de Curso baseada no desenvolvimento da Alvesting, uma plataforma web em que múltiplos modelos de linguagem de larga escala (\textit{Large Language Models}, LLMs) competem entre si por meio de simulações de investimento na B3, a bolsa de valores brasileira. A solução combina, em uma única superfície, um motor de coleta de dados de mercado e notícias em tempo real, um orquestrador de agentes responsável por alimentar cada LLM com o mesmo conjunto de informações, um motor de execução simulada (\textit{paper trading}) que registra ordens fictícias a preços reais de mercado e um painel comparativo de desempenho. A plataforma é organizada em módulos de B3, criptoativos, simulações, conversa com IA, portfólio e notícias, e oferece acesso público com transparência das decisões e respectivas justificativas geradas por cada modelo. Espera-se que a solução contribua para tornar visível, auditável e reproduzível o comportamento de diferentes LLMs no domínio financeiro brasileiro, ainda pouco explorado por iniciativas similares.
\end{abstract-pt}

\begin{abstract-en}
This document presents an undergraduate thesis proposal based on the development of Alvesting, a web platform in which multiple Large Language Models (LLMs) compete against each other through investment simulations on B3, the Brazilian stock exchange. The solution combines, in a single surface, a real-time market data and news ingestion engine, an agent orchestrator that feeds each LLM with the same information set, a paper-trading execution engine that records fictitious orders at real market prices, and a comparative performance dashboard. The platform is organized into modules for B3 equities, crypto assets, simulations, AI chat, portfolio and news, offering public access with full transparency of model decisions and their respective rationales. The expected result is a tool that makes the behavior of different LLMs in the Brazilian financial domain visible, auditable and reproducible, an area still largely unexplored by similar initiatives.
\end{abstract-en}

\begin{pchaves}
apoio à decisão, modelos de linguagem, B3, mercado financeiro, paper trading, agentes autônomos, análise de notícias, plataforma web
\end{pchaves}

\begin{keywords}
decision support, large language models, B3, financial markets, paper trading, autonomous agents, news analysis, web platform
\end{keywords}

\end{frontmatter}

\section{Introdução}

O avanço acelerado dos modelos de linguagem de larga escala nos últimos anos abriu espaço para sua aplicação em domínios antes restritos a sistemas especializados, entre eles o mercado financeiro. Iniciativas recentes têm explorado o uso dessas LLMs como agentes capazes de analisar notícias, dados fundamentalistas e séries históricas para apoiar decisões de investimento, motivando o surgimento de plataformas em que diferentes modelos operam carteiras simuladas e têm seu desempenho comparado de maneira pública. Apesar do interesse crescente nessa abordagem, as iniciativas existentes concentram-se no mercado norte-americano e em criptoativos, deixando o mercado brasileiro praticamente sem cobertura.

No contexto nacional, a tomada de decisão por parte de investidores pessoa física ainda depende fortemente da interpretação manual de notícias, indicadores e cotações dispersos em múltiplas fontes. Esse cenário revela uma limitação informacional importante: ainda que os dados estejam disponíveis em portais e APIs, a integração entre dados de mercado, notícias e análise estruturada permanece pouco acessível ao público geral, e o uso de LLMs como camada de interpretação dessas informações não dispõe de uma vitrine pública adaptada às particularidades da B3.

Nesse contexto, a Alvesting é proposta como uma plataforma web em que múltiplas LLMs --- como GPT, Claude, Gemini, Llama, DeepSeek e Grok --- competem entre si em simulações de investimento na B3. Cada modelo recebe o mesmo conjunto de dados (cotações, fundamentos, indicadores macroeconômicos e notícias nacionais e internacionais) e produz, em cada ciclo de decisão, ordens fictícias acompanhadas de uma justificativa em linguagem natural. A execução é feita em modo \textit{paper trading}, isto é, com dinheiro e ordens fictícios mas a preços reais de mercado, eliminando o risco financeiro sem comprometer o realismo dos resultados. A solução está estruturada como uma aplicação web composta por módulos de B3, criptoativos, simulações, conversa com IA, portfólio e notícias, integrando coleta de dados, orquestração de agentes, execução simulada e visualização comparativa em um único produto.

Dessa forma, o trabalho busca unir desenvolvimento de software e aplicação prática de LLMs em um artefato voltado ao contexto real do mercado financeiro brasileiro, contribuindo tanto para a compreensão pública do comportamento desses modelos no domínio quanto para a aproximação entre tecnologia de ponta e investidor pessoa física.

\section{Problema de Pesquisa}

O problema que motiva este trabalho pode ser resumido da seguinte forma: \textit{como desenvolver uma plataforma capaz de orquestrar múltiplas LLMs em simulações de investimento na B3, alimentando-as com dados reais de mercado e notícias e expondo, de forma transparente e comparável, suas decisões e seus desempenhos ao longo do tempo?}

Na prática, a dificuldade não está apenas em conectar uma LLM a uma corretora ou a um \textit{feed} de cotações, mas em projetar um sistema que (i) trate diferentes provedores de modelos sob uma interface comum, (ii) garanta condições isonômicas de competição entre eles, (iii) consolide dados heterogêneos --- cotações intradiárias, fundamentos por empresa, séries macroeconômicas do Banco Central, notícias em português e em inglês --- em um formato que possa ser entregue às LLMs em cada ciclo de decisão e (iv) registre de maneira auditável cada ordem emitida, a justificativa que a acompanha e o estado da carteira resultante. A esses desafios soma-se a necessidade de que os resultados sejam apresentados ao público de forma compreensível, com indicadores de desempenho como retorno acumulado, comparação com o Ibovespa e evolução temporal do portfólio.

Dessa forma, o desafio do trabalho envolve não apenas integrar APIs e modelos heterogêneos, mas organizá-los em uma arquitetura modular, resiliente e extensível, na qual a adição de uma nova LLM ou de uma nova fonte de dados não exija reescrita do núcleo da plataforma. O problema de pesquisa, portanto, também envolve arquitetura de software, integração e apresentação clara de informações para o usuário final.

\section{Hipótese}

A hipótese deste trabalho é que uma plataforma integrada, capaz de orquestrar múltiplas LLMs e expor publicamente suas decisões em simulações na B3, pode tornar visível e auditável o comportamento desses modelos no domínio financeiro brasileiro. Essa hipótese se sustenta em dois eixos complementares. O primeiro é a \textbf{padronização da arena de competição}: ao garantir que todos os modelos recebam exatamente o mesmo conjunto de dados de entrada (cotações, fundamentos, indicadores e notícias) em cada ciclo de decisão, reduz-se a influência de fatores externos sobre os resultados e torna-se possível atribuir as diferenças de desempenho às características dos próprios modelos. O segundo é a \textbf{transparência de decisões}: ao registrar e exibir publicamente cada ordem emitida e a justificativa em linguagem natural produzida pela LLM, oferece-se ao usuário um material de análise qualitativa que vai além do retorno percentual, permitindo observar não apenas o que cada modelo decide, mas também como ele racionaliza suas escolhas.

Com isso, espera-se que usuários, pesquisadores e entusiastas consigam acompanhar, comparar e debater o desempenho dos diferentes modelos em condições controladas, contribuindo para a compreensão prática das capacidades e limitações das LLMs no mercado financeiro brasileiro.

\section{Objetivos}

\subsection{Objetivo Geral}

Desenvolver a Alvesting como uma plataforma web de competição entre LLMs aplicada ao mercado financeiro brasileiro, oferecendo coleta de dados em tempo real, orquestração de agentes, execução simulada de ordens (\textit{paper trading}) e visualização comparativa do desempenho dos modelos.

\subsection{Objetivos Específicos}

\begin{itemize}
\item Identificar as principais fontes de dados aplicáveis ao contexto brasileiro, incluindo APIs de cotações da B3, séries macroeconômicas do Banco Central, dados fundamentalistas de empresas listadas e feeds de notícias nacionais e internacionais.
\item Projetar uma arquitetura modular capaz de tratar diferentes provedores de LLMs (OpenAI, Anthropic, Google, DeepSeek, xAI, entre outros) sob uma interface comum, utilizando saídas estruturadas (\textit{structured outputs}) para garantir uniformidade nas decisões.
\item Implementar um orquestrador de agentes responsável por construir o contexto de cada ciclo de decisão, distribuí-lo de forma isonômica entre os modelos participantes e coletar suas respostas.
\item Implementar um motor de execução simulada que registre ordens fictícias a preços reais de mercado, atualize a carteira de cada modelo e calcule indicadores de desempenho como retorno acumulado, \textit{drawdown} máximo e comparação com o Ibovespa.
\item Construir uma aplicação web com autenticação, painel comparativo entre modelos, módulos de B3, criptoativos, simulações, conversa com IA, portfólio e notícias, além de páginas públicas que permitam consultar o histórico de decisões e justificativas.
\item Avaliar, em uma janela de operação real, se a plataforma cumpre os requisitos funcionais e não funcionais previstos --- com ênfase em modularidade, resiliência da camada de dados e clareza da apresentação ao usuário.
\end{itemize}

Os objetivos propostos procuram contemplar tanto a construção técnica da solução quanto sua utilidade prática como instrumento de observação e comparação. Dessa forma, o trabalho não se limita à implementação de funcionalidades, mas também busca verificar se o sistema realmente contribui para tornar acessível e auditável o comportamento de LLMs no mercado financeiro brasileiro.

\section{Trabalhos Relacionados}

A seleção de trabalhos relacionados considerou (i) produtos e experimentos públicos em que múltiplas LLMs competem em simulações de investimento e (ii) trabalhos acadêmicos sobre agentes baseados em LLMs aplicados ao mercado financeiro, com ênfase em arquiteturas multiagente e \textit{benchmarks} de desempenho. A Tabela~\ref{tab:relacionados} sumariza os trabalhos selecionados, comparando-os com a Alvesting nas dimensões consideradas relevantes para o domínio do problema: foco específico no mercado brasileiro (B3), competição entre múltiplas LLMs, uso de \textit{paper trading} com dados reais, integração de notícias e modelo de acesso.

\begin{table*}[!ht]
\caption{Comparativo entre Alvesting e trabalhos relacionados.}
\centering
\small
\begin{tblr}{
  width=\textwidth,
  colspec={X[2.4,l] X[1,c] X[1,c] X[1.1,c] X[1.1,c] X[1,c] X[1.1,c]},
  row{1}={font=\bfseries, bg=gray9},
  cells={font=\footnotesize},
  hlines, vlines,
}
Trabalho & Tipo & Foco em B3 & Múltiplas LLMs & \textit{Paper trading} & Notícias & Acesso \\
Rallies AI Arena~\cite{rallies} & Plataforma & \ding{55} & \ding{51} & \ding{51} & \ding{51} & Público \\
Alpha Arena (Nof1)~\cite{nof1} & Experimento & \ding{55} & \ding{51} & \ding{55} & \ding{55} & Público \\
Chen et al. (2025) & Artigo & \ding{55} & \ding{51} & Simulado & \ding{51} & \textit{Open-source} \\
Xiao et al. (2024) & Artigo & \ding{55} & \ding{55} & \ding{51} & \ding{51} & \textit{Open-source} \\
Zhang et al. (2024) & Artigo & \ding{55} & \ding{55} & \ding{51} & \ding{51} & \textit{Open-source} \\
Teles \& Figueiredo (2024) & Artigo & Parcial & \ding{51} & \ding{55} & \ding{51} & \textit{Open-source} \\
\textbf{Alvesting} & Plataforma & \ding{51} & \ding{51} & \ding{51} & \ding{51} & Público \\
\end{tblr}
\label{tab:relacionados}
\end{table*}

A análise da Tabela~\ref{tab:relacionados} evidencia o posicionamento da Alvesting na interseção entre duas categorias de iniciativas que, isoladamente, atendem apenas parte do problema. De um lado, plataformas e experimentos públicos como Rallies AI Arena e Alpha Arena demonstram a viabilidade técnica e o interesse do público em arenas de competição entre LLMs, mas concentram-se respectivamente no mercado de ações dos Estados Unidos e em criptoativos, deixando o mercado brasileiro descoberto. Além disso, a Alpha Arena opera com dinheiro real em uma corretora descentralizada de criptomoedas, o que afasta o experimento da realidade de um investidor pessoa física brasileiro e introduz risco financeiro incompatível com o objetivo educacional pretendido pela Alvesting. De outro lado, trabalhos acadêmicos oferecem fundamentação metodológica para cada um dos eixos centrais da plataforma: Chen et al. (2025) propõem o \textit{benchmark} StockBench para avaliar LLMs como agentes de \textit{trading} em mercados reais; Xiao et al. (2024) organizam diferentes papéis (analistas, pesquisadores e gestores de risco) em uma arquitetura multiagente coesa; Zhang et al. (2024), no \textit{framework} FinAgent publicado na KDD, demonstram a viabilidade de um agente multimodal capaz de integrar dados numéricos, textuais e visuais; e Teles \& Figueiredo (2024), em um estudo conduzido no Brasil, comparam LLMs em tarefas de análise de sentimento aplicadas a notícias de mercado financeiro, sustentando a viabilidade do módulo de Notícias da plataforma proposta. Esses trabalhos, contudo, não constituem plataformas voltadas ao público nem cobrem o domínio brasileiro. A Alvesting posiciona-se justamente nesse cruzamento: combina a transparência e o caráter público das arenas, a fundamentação dos \textit{benchmarks} e \textit{frameworks} acadêmicos e a aderência ao contexto da B3, integrando dados de mercado, notícias em português, múltiplas LLMs e \textit{paper trading} em uma única plataforma. Como contrapartida, o escopo é deliberadamente restrito ao mercado brasileiro, o que se assume simultaneamente como diferencial (especialização no domínio nacional) e como limitação (a solução não generaliza diretamente para outros mercados sem trabalho adicional de integração de dados).

\section{Proposta de Estruturação do Texto}

O texto final do trabalho será organizado em sete capítulos. O Capítulo 1 apresentará a introdução, o problema de pesquisa, a hipótese e os objetivos, definindo o contexto e a motivação do estudo. O Capítulo 2 trará a fundamentação teórica sobre modelos de linguagem, agentes autônomos, mercado financeiro, \textit{paper trading} e sistemas de apoio à decisão. O Capítulo 3 aprofundará os trabalhos relacionados, expandindo a discussão iniciada nesta proposta e posicionando a Alvesting no estado da arte. O Capítulo 4 descreverá a metodologia, os requisitos do sistema e o processo de desenvolvimento adotado. O Capítulo 5 apresentará o desenvolvimento da Alvesting, incluindo arquitetura da aplicação web, camada de coleta de dados, orquestrador de agentes, motor de execução simulada e decisões de implementação. O Capítulo 6 discutirá os resultados obtidos em uma janela de operação real e as limitações da solução. Por fim, o Capítulo 7 reunirá as conclusões e as possibilidades de trabalhos futuros.

Essa organização busca manter uma progressão lógica entre a contextualização do problema, a base conceitual, o posicionamento frente a trabalhos correlatos, o desenvolvimento da solução e a análise dos resultados. Com isso, a estrutura do texto favorece a compreensão do percurso adotado no trabalho e evidencia a relação entre a proposta acadêmica e o artefato desenvolvido.

\section*{Referências}

\begin{thebibliography}{9}

\bibitem{chen2025}
CHEN, Y. et al. \textbf{StockBench: Can LLM Agents Trade Stocks Profitably in Real-World Markets?}. arXiv preprint arXiv:2510.02209, 2025. Disponível em: \url{https://arxiv.org/abs/2510.02209}.

\bibitem{xiao2024}
XIAO, Y. et al. \textbf{TradingAgents: Multi-Agents LLM Financial Trading Framework}. arXiv preprint arXiv:2412.20138, 2024. Disponível em: \url{https://arxiv.org/abs/2412.20138}.

\bibitem{zhang2024}
ZHANG, W. et al. \textbf{A Multimodal Foundation Agent for Financial Trading: Tool-Augmented, Diversified, and Generalist}. In: PROCEEDINGS OF THE 30TH ACM SIGKDD CONFERENCE ON KNOWLEDGE DISCOVERY AND DATA MINING (KDD '24), 2024. DOI: \url{https://doi.org/10.1145/3637528.3672024}.

\bibitem{teles2024}
TELES, L. E. P.; FIGUEIREDO, C. M. S. \textbf{Comparing LLMs for Sentiment Analysis in Financial Market News}. arXiv preprint arXiv:2510.15929, 2024. Disponível em: \url{https://arxiv.org/abs/2510.15929}.

\bibitem{wu2023}
WU, S. et al. \textbf{BloombergGPT: A Large Language Model for Finance}. arXiv preprint arXiv:2303.17564, 2023. Disponível em: \url{https://arxiv.org/abs/2303.17564}.

\bibitem{yang2023}
YANG, H.; LIU, X.-Y.; WANG, C. D. \textbf{FinGPT: Open-Source Financial Large Language Models}. arXiv preprint arXiv:2306.06031, 2023. Disponível em: \url{https://arxiv.org/abs/2306.06031}.

\bibitem{rallies}
LEUNG, K. \textbf{Rallies AI Arena: AI Hedge Fund in Public}. 2026. Disponível em: \url{https://rallies.ai/arena}. Acesso em: 18 maio 2026.

\bibitem{nof1}
NOF1. \textbf{Alpha Arena: LLM Trading Competition Season 1}. 2025. Disponível em: \url{https://nof1.ai/}. Acesso em: 18 maio 2026.

\end{thebibliography}

\end{document}